%%% main.tex --- cons: Cargo-Class Project and Dependency Tooling for Common Lisp
%%%
%%% Portable skeleton: compiles with a stock TeX Live (base `article` + common
%%% packages). Build:  latexmk -pdf main.tex   (or `./build.sh`).

\documentclass[11pt,a4paper]{article}

\usepackage[utf8]{inputenc}
\usepackage[T1]{fontenc}
\usepackage{lmodern}
\usepackage[margin=1in]{geometry}
\usepackage{microtype}
\usepackage{booktabs}
\usepackage{listings}
\usepackage{xcolor}
\usepackage[numbers,sort&compress]{natbib}
\usepackage[hidelinks]{hyperref}

\lstset{basicstyle=\ttfamily\small, columns=fullflexible, keepspaces=true,
  showstringspaces=false, frame=single, framesep=4pt, breaklines=true}

\title{\textbf{cons the magnificent}: Cargo-Class Project\\
       and Dependency Tooling for Common Lisp}
\author{Bob\\ \small\texttt{eternal.recursion@proton.me}}
\date{\today}

\begin{document}
\maketitle

\begin{abstract}
Common Lisp's project and dependency tooling is fragmented across several capable
but disjoint tools---ASDF for system definition and build, Quicklisp for
distribution, ocicl for lockfile-based reproducibility, qlot for pinning, Roswell
for implementation and script management---with no single, modern front-end tying
them together. Contemporary languages solve this with an integrated tool
(\textsc{cargo}, \textsc{npm}, the \textsc{go} tool); Common Lisp has no
equivalent, so every project reinvents its own onboarding scripts. We present
\textsc{cons}, an opinionated front-end that \emph{delegates} to existing engines
and adds developer experience: project scaffolding, integrated build/test/run/repl
tasks driven by a declarative per-project \texttt{cons.lisp} spec (a portable,
GNU-make-free replacement for the per-project Makefile), and a neutral
\emph{dependency-source protocol} with Quicklisp and ocicl backends (and a native
resolver to follow). We describe its design, its place in a Coalton-centered
ecosystem, and a path to lockfile-based reproducibility that could ultimately subsume
the tools it began by wrapping.
\end{abstract}

\section{Introduction}
\label{sec:intro}
% The recurring per-project setup ritual; the absence of a cargo-class tool.

\section{The Fragmentation Problem}
\label{sec:frag}
% ASDF / Quicklisp / ocicl / qlot / Roswell: what each does and doesn't do.

\section{Design: Delegate, then Own}
\label{sec:design}
% A thin front-end; the neutral dependency-source protocol; delegate first,
% native resolver later; the same "pluggable backend" pattern used across the
% ecosystem.

\section{Scaffolding: \texttt{cons init}}
\label{sec:new}
% Templates (lib/cli/web/agent); generating the shared "just works" scaffold.

\section{Build Tasks: the \texttt{cons.lisp} Spec}
\label{sec:tasks}
% A project declares its tasks in a declarative, Lispy `cons.lisp` manifest --- targets
% (build/test/repl/run/serve/dev) that mostly delegate to ASDF/Quicklisp. `cons <target>`
% then runs identically on Linux/macOS/Windows, with no GNU-make dependency: the portable
% replacement for the per-project Makefile. Two execution modes: in-process in cons's warm
% image (a dumped SBCL binary carrying Quicklisp and a large heap, so `cons build` compiles
% Coalton-heavy systems directly), or an isolated subprocess sbcl (`--fresh`, and required
% for save-lisp-and-die targets). Contrast with the Makefile ritual it eliminates.

\section{Dependency Management}
\label{sec:deps}
% Quicklisp and ocicl backends; a native source model + version resolver +
% lockfiles; reproducibility.

\section{Distribution and Self-Hosting}
\label{sec:dist}
% bin/cons via save-lisp-and-die; the bootstrap path; cons managing cons.

\section{The Ecosystem}
\label{sec:ecosystem}
% cons (tooling), hyperion (web), praxeon (agents), a Coalton-first FP library;
% Coalton as a first-class modern Lisp.

\section{Evaluation Plan}
\label{sec:eval}
% Onboarding time vs hand-rolled setup; reproducibility guarantees.

\section{Conclusion}
\label{sec:conclusion}

\bibliographystyle{plainnat}
\bibliography{references}

\end{document}
